\documentclass{ctexart}
\usepackage{template}
\usepackage{esint,extarrows,tikz}
\begin{document}\pagestyle{empty}
\begin{center}\Large
北京大学数学科学学院2022-23高等数学B1期末考试
\end{center}
\begin{problem}[1.(10\songti{分})]
设$\R^3$中平面$x+3y+2z=6$与$x$轴交点为$A$,与$y$轴交点为$B$,与$z$轴交点为$C$.
        \begin{enumerate}[label=\textbf{(\arabic*)},leftmargin=*]
            \item \textbf{(5\songti{分})}\ 求$\triangle ABC$的面积.
            \item \textbf{(5\songti{分})}\ 求过四点$A,B,C,O(0,0,0)$的球面的方程.
        \end{enumerate}
\end{problem}
\begin{solution}
\begin{enumerate}[label=\textbf{(\arabic*)},leftmargin=*]
        \item 分别令$x,y,z$三者中两者为$0$可解得$A(6,0,0),B(0,2,0),C(0,0,3)$.于是
            $$a=|BC|=\sqrt{13},b=|AC|=3\sqrt{5},c=|AB|=2\sqrt{10}$$
            于是$\cos C=\dfrac{a^2+b^2-c^2}{2ab}=\dfrac{18}{6\sqrt{65}}=\dfrac{3}{\sqrt{65}}$,则$\sin C=\sqrt{1-\cos^2C}=\dfrac{2\sqrt{14}}{\sqrt{65}}$.\\
            于是$S_{\triangle ABC}=\dfrac{1}{2}ab\sin C=3\sqrt{14}$.
        \item 设球心$Q(x,y,z)$.根据$|QO|=|QA|=|QB|=|QC|$有
            $$\left\{\begin{array}{l}
                (x-6)^2+y^2+z^2=x^2+y^2+z^2\\
                x^2+(y-2)^2+z^2=x^2+y^2+z^2\\
                x^2+y^2+(z-3)^2=x^2+y^2+z^2
            \end{array}\right.$$
            解得$x=3,y=1,z=\dfrac{3}{2}$.于是球面方程为$\Gamma:(x-3)^2+(y-1)^2+\left(z-\dfrac{3}{2}\right)^2=\dfrac{49}{4}$.
    \end{enumerate}
\end{solution}
\begin{problem}[2.(15\songti{分})]
下面的二元函数的极限存在吗?如果存在,请求出其值;如果不存在,请说明理由.
        \begin{enumerate}[label=\textbf{(\arabic*)},leftmargin=*]
            \item \textbf{(5\songti{分})}\ $\displaystyle\lim_{(x,y)\to(0,0)}\dfrac{24\cos\sqrt{x^2+y^2}-24+12\left(x^2+y^2\right)}{\left(\tan\sqrt{x^2+y^2}\right)^4}$.
            \item \textbf{(5\songti{分})}\ $\displaystyle\lim_{(x,y)\to(0,0)}\left(x+\ln(1+y)\right)\cos\dfrac{1}{x^2+y^2}$.
            \item \textbf{(5\songti{分})}\ $\displaystyle\lim_{(x,y)\to(0,0)}\dfrac{x\sin y}{\sin^2x+\sin^2y}$.
        \end{enumerate}
\end{problem}
\begin{solution}
\begin{enumerate}[label=\textbf{(\arabic*)},leftmargin=*]
        \item 置$z=\sqrt{x^2+y^2}$,于是$(x,y)\to(0,0)$有$z\to 0^+$.于是
            $$\begin{aligned}
                \lim_{(x,y)\to(0,0)}\dfrac{24\cos\sqrt{x^2+y^2}-24+12\left(x^2+y^2\right)}{\left(\tan\sqrt{x^2+y^2}\right)^4}
                &= \lim_{z\to 0^+}\dfrac{24\cos z-24+12z^2}{\tan^4 z} \\
                &= \lim_{z\to 0^+}\dfrac{24\cos z-24+12z^2}{z^4}\cdot\lim_{z\to 0^+}\dfrac{z^4}{\tan^4 z} \\
                &= \lim_{z\to 0^+}\dfrac{-24\sin z+24z}{4z^3}\cdot(1)^4\\
                &= \lim_{z\to 0^+}\dfrac{24-24\cos z}{12z^2} \\
                &= \lim_{z\to 0^+}\dfrac{24\sin z}{24z}\\
                &= 1
            \end{aligned}$$
        \item 我们有$$0\leqslant\left|\left(x+\ln(1+y)\right)\cos\dfrac{1}{x^2+y^2}\right|\leqslant\left|x+\ln(1+y)\right|\leqslant\left|x\right|+\left|\ln(1+y)\right|\leqslant\left|x\right|+2\left|y\right|$$
            而$\displaystyle\lim_{(x,y)\to(0,0)}|x|+|y|=0+0=0$.\\
            其中在$|y|<1/2$时使用$|\ln(1+y)|\leqslant2|y|$.据夹逼定理可知$\displaystyle\lim_{(x,y)\to(0,0)}\left(x+\ln(1+y)\right)\cos\dfrac{1}{x^2+y^2}=0$.
        \item 令$y=kx$,其中$k\neq 0$.于是
            $$\begin{aligned}
                \lim_{(x,y)\to(0,0)}\dfrac{x\sin y}{\sin^2x+\sin^2y}
                &= \lim_{x\to0}\dfrac{x\sin kx}{\sin^2x+\sin^2kx} \\
                &= \lim_{x\to0}\dfrac{k\left(\dfrac{\sin kx}{kx}\right)}{\left(\dfrac{\sin x}{x}\right)^2+k^2\left(\dfrac{\sin kx}{kx}\right)^2} \\
                &= \dfrac{k}{1+k^2}
            \end{aligned}$$
            于是从不同路径接近$(0,0)$时取得的极限不同,因而原极限不存在.
    \end{enumerate}
\end{solution}
\begin{problem}[3.(10\songti{分})]
设$f,g:\R\to\R$都有连续的二阶导数.对于任意$x,y\in\R$,
        定义$h(x,y)=xg\left(\dfrac{y}{x}\right)+f\left(\dfrac{y}{x}\right)$,
        试计算$x^2h_{xx}(x,y)+2xyh_{yx}(x,y)+y^2h_{yy}(x,y)$.
\end{problem}
\begin{solution}
公式在$x\neq0$的区域内有定义.固定$(x,y)$,令$H(s)=h(sx,sy)$,其中$s$在$1$附近,则
\[H(s)=sxg(y/x)+f(y/x)\]
故$H''(1)=0$.另一方面,链式法则给出
\[H''(1)=x^2h_{xx}+2xyh_{yx}+y^2h_{yy}\]
所以所求结果为$0$.现有题干所写的任意$x,y$应理解为$x\neq0$的定义域内.
\end{solution}
\begin{problem}[4.(10\songti{分})]
求$\R^2$中曲线$\e^{xy}+xy+y^2=2$在$(0,1)$处的切线方程.
\end{problem}
\begin{solution}
令$F(x,y)=\e^{xy}+xy+y^2-2$.在$(0,1)$处有$F_x=2,F_y=2$,因此$y'=-F_x/F_y=-1$.
所求切线为$x+y-1=0$.
\end{solution}
\begin{problem}[5.(10\songti{分})]
设三元函数$f(x,y,z)=\left(\dfrac{2x}{z}\right)^{y},z\neq0$.求$f$在点$\left(\dfrac{1}{2},1,1\right)$处下降最快的方向上的单位向量.
\end{problem}
\begin{solution}
由题意$$\dfrac{\p f}{\p x}=\left(\dfrac{2}{z}\right)^yyx^{y-1}\ \ \ \ \ 
    \dfrac{\p f}{\p y}=\left(\dfrac{2x}{z}\right)^y\ln\left(\dfrac{2x}{z}\right)\ \ \ \ \
    \dfrac{\p f}{\p z}=-(2x)^yyz^{-y-1}$$
    将$x=\dfrac12,y=1,z=1$代入可知$\dfrac{\p f}{\p x}=2,\dfrac{\p f}{\p y}=0,\dfrac{\p f}{\p z}=-1$.于是$\tbf{grad}f=\left(2,0,-1\right)$为$f$在该点处的梯度.\\
    于是下降最快的方向与负梯度的方向相同,此方向的单位向量为$\left(-\dfrac{2\sqrt5}{5},0,\dfrac{\sqrt5}{5}\right)$.
\end{solution}
\begin{problem}[6.(10\songti{分})]
求二元函数$f(x,y)=\arctan\dfrac{y}{x}$在点$(2,2)$处的二阶泰勒多项式.
\end{problem}
\begin{solution}
在$(2,2)$处,$f=\pi/4$, $f_x=-1/4$, $f_y=1/4$, $f_{xx}=1/8$, $f_{xy}=0$, $f_{yy}=-1/8$.
因此二阶泰勒多项式为
\[T_2(x,y)=\dfrac\pi4-\dfrac{x-2}{4}+\dfrac{y-2}{4}+\dfrac{(x-2)^2}{16}-\dfrac{(y-2)^2}{16}\]
\end{solution}
\begin{problem}[7.(10\songti{分})]
求函数$f(x)=\left(\sin x\right)^{\frac{2}{3}}+\left(\cos x\right)^{\frac{2}{3}}$在闭区间$\left[-\dfrac{\pi}{2},\dfrac{\pi}{2}\right]$上的最小值,并指明所有最小值点.
\end{problem}
\begin{solution}
令$u=\sin^2x\in[0,1]$,则$f(x)=u^{1/3}+(1-u)^{1/3}$.
由于$t^{1/3}\geqslant t$对$t\in[0,1]$成立,有$f(x)\geqslant1$,等号当且仅当$u=0$或$1$.
所以最小值为$1$,全部最小值点为$-\pi/2,0,\pi/2$.
\end{solution}
\begin{problem}[8.(10\songti{分})]
证明:对于任意给定的$k\in\R$,存在$0$的开邻域$U$和$W$,存在唯一的函数$y=f(x),x\in U,y\in W$满足方程$\e^{kx}+\e^{ky}-2\e^{x+y}=0$.
\end{problem}
\begin{solution}
令$F(x,y)=\e^{kx}+\e^{ky}-2\e^{x+y}$,有$F(0,0)=0$, $F_y(0,0)=k-2$.
当$k\neq2$时,由隐函数定理立即得到原点附近的唯一光滑隐函数.
当$k=2$时,方程等价于$(\e^x-\e^y)^2=0$,故唯一解为$y=x$.
两种情形均满足要求.
\end{solution}
\begin{problem}[9.(15\songti{分})]
设$r$是正实数,$D=\left\{(x,y)\vert\sqrt{x^2+y^2}<r\right\}$,函数$f:D\to\R$满足$f\in C^3(D),f(0,0)=0$,
        $f$在点$(0,0)$处的一阶全微分$\di f(0,0)=0$.$f$在点$(0,0)$处的二阶全微分满足
        $$\di^2f(0,0)=E\left(\Delta x\right)^2+2F\Delta x\Delta y+G\left(\Delta y\right)^2$$
        其中$E,F,G$均为常数.
        \begin{enumerate}[label=\textbf{(\arabic*)},leftmargin=*]
            \item \textbf{(10\songti{分})}\ 证明:存在$D$上的两个函数$a,b:D\to\R$使得$\forall (x,y)\in D$有$$f(x,y)=xa(x,y)+yb(x,y),a(0,0)=b(0,0)=0$$
            \item \textbf{(5\songti{分})}\ 若$E>0,EG-F^2<0$,则在$\R^3$中点$(0,0,0)$的充分小邻域中,曲面$z=f(x,y)$充分近似于哪一类二次曲面?画出此类二次曲面的草图.
                从此类二次曲面的几何形状判断是否存在$\R^2$中点$(0,0)$的充分小邻域$D_1$,存在$D_1$上的一一对应的$C^1$变量变换$x=x(u,v),y=y(u,v)$使得
                $$f(x(u,v),y(u,v))=u^2-v^2$$
        \end{enumerate}
\end{problem}
\begin{solution}
\begin{enumerate}[label=\tbf{(\arabic*)},leftmargin=*]
\item 定义
\[a(x,y)=\int_0^1f_x(tx,ty)\di t,\qquad b(x,y)=\int_0^1f_y(tx,ty)\di t\]
圆盘$D$关于原点星形,故积分有定义.由微积分基本定理,
\[xa(x,y)+yb(x,y)=\int_0^1\dfrac{\di}{\di t}f(tx,ty)\di t=f(x,y)-f(0,0)=f(x,y)\]
而$a(0,0)=f_x(0,0)=0$, $b(0,0)=f_y(0,0)=0$.
\item 二阶泰勒公式给出$f(x,y)=\frac12(Ex^2+2Fxy+Gy^2)+o(x^2+y^2)$.
由于$EG-F^2<0$,二次型一正一负,相应二次曲面是双曲抛物面,外形为马鞍面.
\begin{center}
\begin{tikzpicture}[scale=1.25]
\foreach \v in {-1,-.5,0,.5,1}{
  \draw[gray] plot[domain=-1:1,samples=30] (\x-\v,{.4*(\x+\v)+.5*(\x*\x-\v*\v)});
  \draw[gray] plot[domain=-1:1,samples=30] (\v-\x,{.4*(\v+\x)+.5*(\v*\v-\x*\x)});
}
\draw[->] (0,0)--(1.6,.64) node[right] {$u$};
\draw[->] (0,0)--(-1.6,.64) node[left] {$v$};
\draw[->] (0,0)--(0,1.6) node[above] {$z$};
\end{tikzpicture}
\end{center}
精确等式一般不能仅靠线性旋转伸缩得到.为证明存在所要求的$C^1$变换,由带积分余项的Taylor公式,写成
\[f(x,y)=A(x,y)x^2+2B(x,y)xy+C(x,y)y^2\]
其中$A=\int_0^1(1-t)f_{xx}(tx,ty)\di t$, $B,C$分别以$f_{xy},f_{yy}$代替$f_{xx}$定义.
它们为$C^1$函数,且原点处$A=E/2>0$, $AC-B^2=(EG-F^2)/4<0$.
在足够小的邻域内定义
\[u=\sqrt A\left(x+\dfrac BAy\right),\qquad
v=y\sqrt{B^2/A-C}\]
则精确满足$f=u^2-v^2$.此映射在原点的Jacobian为$\sqrt{B(0)^2-A(0)C(0)}>0$.
由逆函数定理,存在局部一一对应的$C^1$逆变换$x=x(u,v),y=y(u,v)$满足题意.
\end{enumerate}
\end{solution}
\end{document}
